% © 2026 Ing. David Jaroš — CC BY-NC-ND 4.0
%
% This work is licensed under a Creative Commons Attribution-NonCommercial-NoDerivatives
% 4.0 International License (CC BY-NC-ND 4.0).
%
% See LICENSE.md for full license text.
%
% File: research_tracks/T3_ALPHA/modular_symmetry_rpsi_EN.tex
% Purpose: Systematic probe — modular S-symmetry of S[Theta] as mechanism
%          for fixing R_psi = 1 (self-dual point).
% Status:  Path 1 (polynomial): [L0]+[MC].
%          Path 2 (Z[tau]=eta^{-12}): [L1 conditional on one-loop + STD].
%          R_psi=1: derived via two independent paths.
%          Alpha NOT DERIVED. G137-B NOT CLOSED.
% Date:    2026-06-11 (updated: Z[tau] chain added)

\ifdefined\INCLUDEMODE
\else
\documentclass[12pt,a4paper]{article}
\usepackage[a4paper,margin=2.5cm]{geometry}
\usepackage{amsmath,amssymb,amsthm}
\usepackage{mathtools}
\usepackage{hyperref}
\usepackage{mdframed}
\usepackage{booktabs}
\usepackage{xcolor}
\usepackage{microtype}

\hypersetup{
  colorlinks=true,
  linkcolor=blue!60!black,
  citecolor=green!50!black,
  urlcolor=blue!60!black
}

\newtheorem{theorem}{Theorem}[section]
\newtheorem{lemma}[theorem]{Lemma}
\newtheorem{proposition}[theorem]{Proposition}
\newtheorem{corollary}[theorem]{Corollary}
\theoremstyle{definition}
\newtheorem{definition}[theorem]{Definition}
\theoremstyle{remark}
\newtheorem{remark}[theorem]{Remark}

\newmdenv[backgroundcolor=yellow!20, linecolor=orange, linewidth=1.5pt]{gapbox}
\newmdenv[backgroundcolor=green!15, linecolor=green!60!black, linewidth=1.5pt]{resultbox}
\newmdenv[backgroundcolor=red!10, linecolor=red!60, linewidth=1.5pt]{deadendbox}
\newmdenv[backgroundcolor=blue!8, linecolor=blue!50, linewidth=1.5pt]{insightbox}

\newcommand{\PROVED}{\textbf{[Proved]}}
\newcommand{\OPEN}{\textbf{[Open]}}
\newcommand{\Lzero}{\textbf{[L0]}}
\newcommand{\Lone}{\textbf{[L1]}}
\newcommand{\Ltwo}{\textbf{[L2]}}
\newcommand{\MC}{\textbf{[MC]}}
\newcommand{\NUM}{\textbf{[NUM]}}
\newcommand{\OBS}{\textbf{[OBS]}}
\newcommand{\COND}{\textbf{[L1 conditional]}}
\newcommand{\STD}{\textbf{[STD]}}
\newcommand{\NODERIVED}{\textbf{NOT DERIVED}}

\newcommand{\CC}{\mathbb{C}}
\newcommand{\HH}{\mathbb{H}}
\newcommand{\RR}{\mathbb{R}}
\newcommand{\ZZ}{\mathbb{Z}}

\title{\textbf{Modular Symmetry of $S[\Theta]$ as a Mechanism for Fixing $R_\psi$}\\[6pt]
  \large Systematic probe: $\tau \to -1/\tau$ and the self-dual point $R_\psi = 1$}
\author{Ing.\ David Jaroš}
\date{2026-06-11}

\begin{document}
\maketitle

\begin{abstract}
We investigate the fixing of $R_\psi = 1$ (self-dual point $\tau_E = i$)
via two independent paths.

\textbf{Path 1 (Casimir polynomial)}: If the effective action $S_{\mathrm{eff}}$
is S-invariant ($R_\psi \to 1/R_\psi$), stationarity yields
$R_\psi^6 + R_\psi^4 + R_\psi^2 = 3$ — the unique positive solution is $R_\psi = 1$.
$N_{\mathrm{eff}}$ cancels; the result is an algebraic property of S-symmetry.
Status: \textbf{[L0]} (polynomial) $+$ \textbf{[MC]} (S-invariance of $S_{\mathrm{eff}}$).

\textbf{Path 2 (partition function)}: From the one-loop partition function
$Z[\tau_E] = \eta(\tau_E)^{-12}$ (\textbf{[L1]}: $\det P_{\mathrm{eff}} = \eta^2$
\textbf{[STD]} $+$ $N_{\mathrm{eff}} = 12$ \textbf{[L1]}), the transformation
property $\eta(-1/\tau) = \sqrt{-i\tau}\,\eta(\tau)$ makes $\tau_E = i$ a
fixed point of the S-transformation. $\psi$ is compact from the SS boundary
condition \textbf{[L1]} and from the periodicity of $Z[\tau]$ \textbf{[L0]}.
Status: \textbf{[L1 conditional on one-loop approximation and STD result]}.

The agreement of both paths is structural evidence. Implications:
scale closure \textbf{[MC]}; $\tau = i$ algebraically preferred for
alpha \textbf{[OBS]}.
\textbf{Alpha NOT DERIVED.}
\end{abstract}

\tableofcontents

%──────────────────────────────────────────────────────────────────────────────
\section{Context and Motivation}
\label{sec:context}
%──────────────────────────────────────────────────────────────────────────────

Three open problems in UBT share a common root: the radius $R_\psi$ of the
$\psi$-circle is not derived from $S[\Theta]$ without phenomenological input.

\begin{enumerate}
  \item \textbf{Alpha (G137-B)}: $B_{\mathrm{phenom}}$ depends on the
    self-dual point $\tau = i$, which occurs at $R_\psi = 1$. Why is $\tau = i$
    selected? Not derived.
  \item \textbf{Scale closure (EW-1b)}: $M_{\mathrm{GUT}} = 1/R_\psi^*$;
    the stationarity equation gives $R_\psi^*$ conditionally on $T_{\mathrm{kin}}$,
    which is not derived from $S[\Theta]$.
  \item \textbf{C2-iv}: $\psi$-winding and the colour fiber are decoupled.
    A connection might arise from the geometry of the $\psi$-circle at a
    special value of $R_\psi$.
\end{enumerate}

\textbf{Hypothesis}: $S[\Theta]$ has modular symmetry under $\tau \to -1/\tau$
(inversion of complex time). If so, the self-dual point $\tau = i$ is a
\emph{fixed point} of this symmetry — $R_\psi = 1$ would be derived from the
symmetry of $S[\Theta]$, not chosen by hand.

\begin{gapbox}
\textbf{Scope of the probe.}
This probe tests \emph{modular symmetry of the effective action $S_{\mathrm{eff}}$}
as a mechanism for fixing $R_\psi$. We do not attempt to derive $B$ directly
from modular forms (Route A3, NO-GO). This is a new, previously untested route.
\end{gapbox}

%──────────────────────────────────────────────────────────────────────────────
\section{Modular Transformation in the UBT Context}
\label{sec:setup}
%──────────────────────────────────────────────────────────────────────────────

\subsection*{Complex time and the $\psi$-circle}

Complex time $\tau = t + i\psi \in \CC$. The $\psi$-circle $S^1_\psi$ has
period $2\pi R_\psi$. The Euclidean modular parameter:
\[
  \tau_E = iR_\psi \in \mathbb{H},
  \quad t = 0,\; \psi \sim \psi + 2\pi R_\psi.
\]

Standard modular transformations on the upper half-plane $\mathbb{H}$:
\[
  T:\; \tau_E \to \tau_E + 1, \qquad
  S:\; \tau_E \to -\frac{1}{\tau_E}.
\]
The group $\mathrm{SL}(2,\ZZ)$ is generated by $T$ and $S$.

\begin{definition}[Modular S-transformation on UBT complex time]
\label{def:modular_S}
The S-transformation on the $\psi$-circle:
\[
  \tau_E \;\longmapsto\; -\frac{1}{\tau_E},
  \qquad
  iR_\psi \;\longmapsto\; \frac{i}{R_\psi},
\]
equivalently $R_\psi \to 1/R_\psi$ — T-duality on the $\psi$-circle
\Lone~(standard T-duality~\cite[§8]{polchinski_v1}).

The self-dual point of the S-transformation: $S(i) = -1/i = i$, so $R_\psi = 1$
is a \emph{fixed point} of the S-transformation.
\end{definition}

%──────────────────────────────────────────────────────────────────────────────
\section{Test of Modular Invariance of the Kinetic Term}
\label{sec:kinetic}
%──────────────────────────────────────────────────────────────────────────────

The kinetic term of $S[\Theta]$ (from \texttt{appendix\_AA\_theta\_action.tex}):
\[
  S_{\mathrm{kin}} = \int \mathrm{Tr}\!\left[(D_\mu\Theta)^\dagger(D^\mu\Theta)\right]
  d^4x\, d\psi.
\]
Integration over $d\psi$ with period $2\pi R_\psi$.

\begin{proposition}[Transformation of the kinetic term under $R_\psi \to 1/R_\psi$]
\label{prop:kin_transform}
Under the S-transformation $R_\psi \to 1/R_\psi$:
\begin{itemize}
  \item Period: $2\pi R_\psi \to 2\pi/R_\psi$.
  \item KK spectrum: $m_n = n/R_\psi \to n R_\psi$ (inversion of masses).
  \item Volume: $\int d\psi = 2\pi R_\psi \to 2\pi/R_\psi$.
\end{itemize}
Therefore:
\[
  S_{\mathrm{kin}}[\Theta,\, R_\psi]
  \;\xrightarrow{\;R_\psi \,\to\, 1/R_\psi\;}
  \frac{1}{R_\psi^2}\, S_{\mathrm{kin}}[\widetilde\Theta,\, 1/R_\psi],
\]
where $\widetilde\Theta$ is the S-dual field. The factor $R_\psi^{-2}$ arises
from dimensional analysis: $[\Theta] = M^1$ and the field rescales as
$\widetilde\Theta = R_\psi \cdot \Theta$ \MC.
$S_{\mathrm{kin}}$ is S-invariant if $\widetilde\Theta = R_\psi \cdot \Theta$.
\end{proposition}

%──────────────────────────────────────────────────────────────────────────────
\section{$\psi$-Sector Partition Function and First-Principles Fixing of $R_\psi = 1$}
\label{sec:Z_chain}
%──────────────────────────────────────────────────────────────────────────────

This section provides an \emph{independent} derivation of $R_\psi = 1$, independent
of the Casimir potential and S-invariance of $S_{\mathrm{eff}}$ from
\S\ref{sec:casimir}. The derivation follows from the algebraic structure of the
UBT $\psi$-sector partition function.

%──────────────────────────────────────────────────────────────────────────────
\subsection{Compactness of $\psi$ from the Scherk-Schwarz condition}

\begin{lemma}[Compactness of $\psi$ from the SS boundary condition — \Lone]
\label{lem:SS_BC}
From \texttt{su2\_twist\_neff12.tex}~\cite{su2_twist},
\S\,``Twisted boundary condition on $S^1_\psi$'':
\[
  \Theta(\psi + 2\pi R_\psi) = \sigma_3\,\Theta(\psi)\,\sigma_3,
  \qquad \sigma_3 = \mathrm{diag}(+1,-1).
\]
This condition gives $\psi$ a natural period $2\pi R_\psi$:
the off-diagonal components $b, c$ are anti-periodic
($m_n = (n+\tfrac12)/R_\psi$), and the diagonal components $a, d$ are periodic
($m_n = n/R_\psi$).
\textbf{The compactness of $\psi$ follows from the symmetry of the field
$\Theta$, not from a geometric axiom.} \Lone
\end{lemma}

%──────────────────────────────────────────────────────────────────────────────
\subsection{UBT $\psi$-sector partition function}

\begin{lemma}[One-loop partition function of a scalar field on $S^1$ — \STD]
\label{lem:Z_scalar}
For a single periodic real scalar field on $S^1$ with modular parameter
$\tau_E = iR_\psi$, the one-loop functional determinant
is~\cite[§7]{polchinski_v1}:
\[
  \det P_{\mathrm{eff}} = \eta(\tau_E)^2,
\]
where $\eta(\tau) = q^{1/24}\prod_{n=1}^\infty(1-q^n)$, $q = e^{2\pi i\tau}$.
The one-loop partition function is therefore:
\[
  -\ln Z_{1\text{-loop}} = \tfrac12 \ln \det P_{\mathrm{eff}}
  = \ln \eta(\tau_E).
\]
\end{lemma}

\begin{proposition}[Total UBT $\psi$-sector partition function — \Lone]
\label{prop:Z_total}
From Lemma~\ref{lem:Z_scalar} and $N_{\mathrm{eff}} = 12$
\Lone~\cite{su2_twist}:
\[
  \boxed{Z[\tau_E] \;=\; \eta(\tau_E)^{-N_{\mathrm{eff}}} \;=\; \eta(\tau_E)^{-12}.}
\]
\textbf{Argument}: The one-loop contribution of $N_{\mathrm{eff}}$ effective
modes (with helicity weights from the SS twist condition) to the partition
function is:
\[
  -\ln Z = \sum_{k=1}^{N_{\mathrm{eff}}} \ln \eta(\tau_E)
  = N_{\mathrm{eff}} \ln \eta(\tau_E).
\]
$N_{\mathrm{eff}} = 3 \times 2 \times 2 = 12$ incorporates three compact phase
directions $(\psi, \chi, \xi)$, two charged off-diagonal components $(b, c)$,
and two Weyl helicity weights from the half-integer tower —
all \Lone~\cite{su2_twist}.

\textbf{Note}: $Z[\tau_E] = \eta(\tau_E)^{-12}$ is conditional on the
one-loop approximation (higher loops correct but do not shift the fixed point)
and on $\det P_{\mathrm{eff}} = \eta^2$ \STD.
\end{proposition}

%──────────────────────────────────────────────────────────────────────────────
\subsection{Effective compactification from periodicity of $Z[\tau]$}

\begin{proposition}[Effective compactification of $\psi$ from periodicity
of $Z[\tau]$ — \Lzero]
\label{prop:psi_compact_from_Z}
From the transformation property of $\eta$~\cite[§9.2]{apostol_modular}:
\[
  \eta(\tau_E + 1) = e^{i\pi/12}\,\eta(\tau_E).
\]
Therefore for $Z[\tau_E] = \eta(\tau_E)^{-12}$:
\[
  Z[\tau_E + 1] = e^{-i\pi}\,Z[\tau_E] = -Z[\tau_E].
\]
The physical amplitude $|Z[\tau_E + 1]|^2 = |Z[\tau_E]|^2$ — periodic
in $\tau_E \to \tau_E + 1$. This periodicity induces an effective identification
$\tau_E \sim \tau_E + 1$ as a symmetry of physical observables. \Lzero

\textbf{Consequence}: $\psi$ is effectively compact both from the SS condition
(Lemma~\ref{lem:SS_BC}, \Lone) and from the periodicity of $Z[\tau]$ (\Lzero).
Both mechanisms are consistent and independent.
\end{proposition}

%──────────────────────────────────────────────────────────────────────────────
\subsection{Modular transformation of $Z[\tau]$ and fixing $R_\psi = 1$}

\begin{proposition}[Modular transformation of $Z[\tau]$ and stationarity
at $\tau_E = i$ — \Lzero]
\label{prop:Z_modular}
From the transformation property of $\eta$~\cite[§9.2]{apostol_modular}:
\[
  \eta\!\left(-\tfrac{1}{\tau_E}\right) = \sqrt{-i\tau_E}\;\eta(\tau_E).
\]
Therefore:
\[
  Z\!\left[-\tfrac{1}{\tau_E}\right]
  = (-i\tau_E)^{-N_{\mathrm{eff}}/2}\,Z[\tau_E].
\]
With $S[\Theta] = -\ln Z[\tau_E]$:
\[
  S\!\left[-\tfrac{1}{\tau_E}\right]
  = S[\tau_E] + \tfrac{N_{\mathrm{eff}}}{2}\ln(-i\tau_E).
\]
For $\tau_E = iR_\psi$: $\,-i\tau_E = R_\psi > 0$, so
$\ln(-i\tau_E) = \ln R_\psi$. At the self-dual point $R_\psi = 1$:
\[
  \ln(-i\cdot i) = \ln 1 = 0
  \;\;\Longrightarrow\;\;
  \boxed{S[\tau_E = i] \text{ is a stationary point of the modular flow.}}
\]
\Lzero\ — algebraic identity.
\end{proposition}

\begin{theorem}[Fixing $R_\psi = 1$ from the symmetry of $Z[\tau]$ — \Lone]
\label{thm:Rpsi_from_Z}
Under the assumptions of Proposition~\ref{prop:Z_total} (\Lone):
\begin{enumerate}
  \item $\tau_E = i$ ($R_\psi = 1$) is a fixed point of the S-transformation
    $\tau_E \to -1/\tau_E$ on $Z[\tau_E]$: the factor
    $(-i\tau_E)^{-N_{\mathrm{eff}}/2} = 1$ for $\tau_E = i$. \Lzero
  \item $\tau_E = i$ is the \emph{only} point on $\mathbb{H}$ where
    $Z[\tau_E]$ is invariant under the S-transformation. \Lzero\
    (The fixed points of $\mathrm{SL}(2,\ZZ)$ on $\mathbb{H}$ are $\tau = i$
    and $\tau = e^{2\pi i/3}$; for $\tau_E = iR_\psi$ with $R_\psi > 0$
    only $\tau_E = i$ lies on the imaginary axis.)
  \item Combined with Theorem~\ref{thm:poly_fixation} (polynomial — two
    independent paths to $R_\psi = 1$): agreement confirms robustness.
\end{enumerate}
\textbf{Conclusion}: $R_\psi = 1$ is derived from the symmetry of the
partition function $Z[\tau_E] = \eta(\tau_E)^{-12}$.
Status: \Lone\ conditional on the one-loop approximation and $\det P = \eta^2$
\STD; unconditionally \Lone\ upon citation of the standard result.
\end{theorem}

\begin{remark}[Two independent paths to $R_\psi = 1$]
\label{rem:two_paths}
\begin{center}
\begin{tabular}{lll}
  \toprule
  Path & Mechanism & Status \\
  \midrule
  Polynomial (\S\ref{sec:stationary}) &
    Casimir dynamics $+$ S-invariance of $S_{\mathrm{eff}}$ &
    \Lzero\ $+$ \MC \\
  $Z[\tau]$ (\S\ref{sec:Z_chain}) &
    Modular symmetry of partition function &
    \Lone\ (cond.) \\
  \bottomrule
\end{tabular}
\end{center}
The agreement of two independent computations is structural evidence of
robustness: $R_\psi = 1$ is the natural scale of the UBT $\psi$-sector.
\end{remark}

%──────────────────────────────────────────────────────────────────────────────
\section{Casimir Term and S-Invariance Condition for $S_{\mathrm{eff}}$}
\label{sec:casimir}
%──────────────────────────────────────────────────────────────────────────────

\subsection*{Casimir potential}

From \texttt{research\_tracks/EW/rpsi\_from\_action.tex} and the standard
Casimir energy on $T^3$ with $d = 3$ \Lone~\cite{ambjorn_wolfram_1983}:
\[
  V_{\mathrm{Casimir}}(R_\psi) = \frac{C_{\mathrm{total}}}{R_\psi^3},
  \qquad
  C_{\mathrm{total}} = N_{\mathrm{eff}} \cdot C_{\mathrm{unit}},
  \qquad
  C_{\mathrm{unit}} = \frac{3\,\zeta(3)}{128\pi^2}.
\]
($C_{\mathrm{total}} > 0$ corresponds to standard Casimir repulsion away from $R_\psi = 0$.)

\subsection*{Effective action}

\[
  S_{\mathrm{eff}}(R_\psi,\, T_{\mathrm{kin}}) = 2\pi R_\psi \cdot T_{\mathrm{kin}}
  + \frac{C_{\mathrm{total}}}{R_\psi^3}.
\]

\subsection*{S-invariance of $S_{\mathrm{eff}}$}

\begin{lemma}[S-invariance condition for $S_{\mathrm{eff}}$]
\label{lem:S_inv_condition}
$S_{\mathrm{eff}}[R_\psi,\, T_{\mathrm{kin}}] = S_{\mathrm{eff}}[1/R_\psi,\, T_{\mathrm{kin}}]$
if and only if:
\[
  2\pi T_{\mathrm{kin}}\!\left(R_\psi - \frac{1}{R_\psi}\right)
  = V_{\mathrm{Casimir}}\!\left(\frac{1}{R_\psi}\right) - V_{\mathrm{Casimir}}(R_\psi)
  = C_{\mathrm{total}}\!\left(R_\psi^3 - \frac{1}{R_\psi^3}\right).
\]
For $R_\psi \neq 1$:
\begin{equation}
  \label{eq:T_kin_S_inv}
  T_{\mathrm{kin}}(R_\psi) = \frac{C_{\mathrm{total}}}{2\pi}
  \cdot \frac{R_\psi^3 - R_\psi^{-3}}{R_\psi - R_\psi^{-1}}
  = \frac{C_{\mathrm{total}}}{2\pi}
  \left(R_\psi^2 + 1 + \frac{1}{R_\psi^2}\right).
\end{equation}
L'H\^{o}pital limit for $R_\psi \to 1$:
$T_{\mathrm{kin}}(1) = 3C_{\mathrm{total}}/(2\pi)$.
\end{lemma}

\begin{proof}
Direct computation:
$S_{\mathrm{eff}}[R] - S_{\mathrm{eff}}[1/R]
= 2\pi T_{\mathrm{kin}}(R - 1/R) + C_{\mathrm{total}}(1/R^3 - R^3) = 0$
implies \eqref{eq:T_kin_S_inv}.
Factorization $R^3 - 1/R^3 = (R - 1/R)(R^2 + 1 + 1/R^2)$ \Lzero.
\end{proof}

%──────────────────────────────────────────────────────────────────────────────
\section{S-Invariance Condition, Stationarity, and the Polynomial}
\label{sec:stationary}
%──────────────────────────────────────────────────────────────────────────────

\subsection*{Stationarity equation}

Stationarity $dS_{\mathrm{eff}}/dR_\psi = 0$:
\[
  \frac{d}{dR_\psi}\!\left[2\pi R_\psi T_{\mathrm{kin}}
  + \frac{C_{\mathrm{total}}}{R_\psi^3}\right] = 0
  \;\;\Longrightarrow\;\;
  2\pi T_{\mathrm{kin}} = \frac{3C_{\mathrm{total}}}{R_\psi^4}.
\]

\subsection*{Key substitution}

Substituting $T_{\mathrm{kin}}$ from the S-invariance condition~\eqref{eq:T_kin_S_inv}:
\[
  \frac{C_{\mathrm{total}}}{2\pi}\cdot 2\pi \cdot
  \left(R_\psi^2 + 1 + \frac{1}{R_\psi^2}\right)
  = \frac{3C_{\mathrm{total}}}{R_\psi^4}.
\]
$C_{\mathrm{total}}$ cancels (independent of $R_\psi$):
\[
  R_\psi^2 + 1 + \frac{1}{R_\psi^2} = \frac{3}{R_\psi^4}.
\]
Multiplying by $R_\psi^4$:

\begin{equation}
  \label{eq:poly}
  \boxed{R_\psi^6 + R_\psi^4 + R_\psi^2 = 3.}
\end{equation}

\begin{theorem}[Polynomial fixing of $R_\psi$ from S-invariance]
\label{thm:poly_fixation}
Conditional on S-invariance of $S_{\mathrm{eff}}$ and the Casimir form
$V = C_{\mathrm{total}}/R_\psi^3$, $R_\psi$ satisfies the polynomial equation
\eqref{eq:poly}. This equation has the unique positive real solution $R_\psi = 1$.

Key fact: $N_{\mathrm{eff}}$ \emph{cancels} (it appears only in $C_{\mathrm{total}}$,
which divides out) — fixing $R_\psi = 1$ is an algebraic property of S-symmetry,
independent of $N_{\mathrm{eff}}$.
\end{theorem}

\begin{proof}
\emph{$R_\psi = 1$ is a solution}: $1 + 1 + 1 = 3$ \Lzero.
\emph{Uniqueness}: the substitution $u = R_\psi^2 > 0$ gives $u^3 + u^2 + u = 3$,
i.e.\ $f(u) = u^3 + u^2 + u - 3$. $f(1) = 0$, $f'(u) = 3u^2 + 2u + 1 > 0$
for all $u$ (negative discriminant) — $f$ is \emph{strictly increasing} on
$(0,\infty)$, so $u = 1$ is the unique positive root. \Lzero.
\emph{$N_{\mathrm{eff}}$ cancels}: $C_{\mathrm{total}} = N_{\mathrm{eff}} C_{\mathrm{unit}}$
appears on both sides of the stationarity equation and divides out.
Numerically verified for $N_{\mathrm{eff}} \in \{1, 3, 6, 12, 24, 48\}$
in \texttt{tools/modular\_rpsi\_check.py}.
\end{proof}

\begin{corollary}[$R_\psi = 1$ is a minimum]
\label{cor:minimum}
$R_\psi = 1$ is a \emph{minimum} (not a maximum) of $S_{\mathrm{eff}}$ with
$T_{\mathrm{kin}}$ from S-invariance. Second derivative:
\[
  \frac{d^2 S_{\mathrm{eff}}}{dR_\psi^2}\bigg|_{R_\psi=1} = 20\, C_{\mathrm{total}} > 0.
\]
\emph{Calculation}: $d^2T_{\mathrm{kin}}/dR^2|_{R=1} = 4C_{\mathrm{total}}/\pi$ and
$dT_{\mathrm{kin}}/dR|_{R=1} = 0$ \Lzero.
Thus $d^2 S/dR^2 = 2\pi \cdot (0 + 4C/\pi) + 12C = 8C + 12C = 20C > 0$.
Numerically: $d^2S/dR^2|_{R=1} = 20 \times 0.03425 \approx 0.685 > 0$. \NUM
\end{corollary}

\subsection*{Resolution of the apparent $N_{\mathrm{eff}} = 1$ inconsistency}

\begin{remark}[Notational inconsistency in the naive analysis]
\label{rem:notation}
A naive analysis that conflates $C_{\mathrm{unit}}$ (per mode) and $C_{\mathrm{total}}$
(total) yields:
\[
  R^6 + R^4 + R^2 = 3N_{\mathrm{eff}} \;\Rightarrow\; N_{\mathrm{eff}} = 1
  \;\text{(for $R=1$)}.
\]
This is a \emph{notational artefact}: if Lemma~\ref{lem:S_inv_condition}
uses $C_{\mathrm{unit}}$ (per mode) but the stationarity equation uses
$3N_{\mathrm{eff}}C_{\mathrm{unit}}/R^4$, then $N_{\mathrm{eff}}$ does not cancel.
The correct analysis with consistent $C_{\mathrm{total}}$ yields
equation~\eqref{eq:poly}, in which $N_{\mathrm{eff}}$ is \emph{absent} —
fixing $R=1$ holds for any $N_{\mathrm{eff}}$, including $N_{\mathrm{eff}} = 12$.
\end{remark}

%──────────────────────────────────────────────────────────────────────────────
\section{Gauge Term Contribution}
\label{sec:gauge}
%──────────────────────────────────────────────────────────────────────────────

$S[\Theta]$ contains the gauge term
$S_{\mathrm{gauge}} = \int \mathrm{Tr}[F_{\mu\nu}F^{\mu\nu}]\,d^4x\,d\psi$.
Under the S-transformation $R_\psi \to 1/R_\psi$: the gauge coupling
$g \to g/R_\psi$ (KK rescaling)~\cite[§12]{polchinski_v1}.

\begin{lemma}[Effective $T_{\mathrm{kin}}^{\mathrm{eff}}$ with gauge contribution]
\label{lem:Tkin_eff}
Including the gauge term as a contribution to the effective kinetic term:
\[
  T_{\mathrm{kin}}^{\mathrm{eff}} = T_{\mathrm{kin}} + \Delta T_{\mathrm{gauge}},
  \quad
  \Delta T_{\mathrm{gauge}} \sim N_{\mathrm{colour}}\,\frac{g^2 M_{\mathrm{KK}}^4}{16\pi^2},
\]
where $N_{\mathrm{colour}} = 8\,(\text{gluons}) + 3\,(\text{weak}) + 1\,(\text{hypercharge}) = 12$.

\textbf{Observation}: $N_{\mathrm{colour}} = N_{\mathrm{eff}} = 12$ \MC.
This equality may not be coincidental: both counts reflect the number of
physical bosons of $\mathrm{SU}(3) \times \mathrm{SU}(2)_L \times \mathrm{U}(1)_Y$
in UBT with $N_{\mathrm{eff}} = 12$ [L1 conditional on gauge structure from UBT].
\end{lemma}

\begin{remark}[Gauge contribution and the polynomial]
\label{rem:gauge_poly}
Adding $\Delta T_{\mathrm{gauge}}$ modifies $T_{\mathrm{kin}}$ in
Lemma~\ref{lem:S_inv_condition}. If $\Delta T_{\mathrm{gauge}}$ is itself
S-invariant (i.e.\ $\Delta T_{\mathrm{gauge}}(R) = \Delta T_{\mathrm{gauge}}(1/R)/R^2$
after appropriate field rescaling), then Theorem~\ref{thm:poly_fixation} remains
valid with $C_{\mathrm{total}} \to C_{\mathrm{total}} + C_{\mathrm{gauge}}$.
The condition $R = 1$ and polynomial~\eqref{eq:poly} are robust under such
S-invariant corrections. Derivation of $C_{\mathrm{gauge}}$ from $S[\Theta]$: \OPEN.
\end{remark}

%──────────────────────────────────────────────────────────────────────────────
\section{Overall Verdict: S-Invariance of $S[\Theta]$?}
\label{sec:verdict}
%──────────────────────────────────────────────────────────────────────────────

\begin{insightbox}
\textbf{Result C — Conditional S-invariance}

$S_{\mathrm{eff}}$ is S-invariant \emph{conditional on}:
\begin{enumerate}
  \item $S[\Theta]$ is invariant under $\tau_E \to -1/\tau_E$ (not yet derived
    from the symmetry of $S[\Theta]$; requires derivation from first principles).
  \item The Casimir potential $V_{\mathrm{Casimir}} = C_{\mathrm{total}}/R_\psi^3$
    is derivable from $S[\Theta]$ (standard result for $T^3$,
    \cite{ambjorn_wolfram_1983}, but explicit derivation from $S[\Theta]$: \OPEN).
  \item $T_{\mathrm{kin}}$ is determined by the S-invariance condition
    (self-consistency; independent value of $T_{\mathrm{kin}}$ from $S[\Theta]$: \OPEN).
\end{enumerate}

Conditionally on (1)--(3):
\[
  R_\psi = 1 \text{ (self-dual point)}
  \;\text{ is the minimum of } S_{\mathrm{eff}},
  \quad
  T_{\mathrm{kin}}(R_\psi{=}1) = \frac{3C_{\mathrm{total}}}{2\pi}.
\]

\textbf{Key algebraic fact} (independent of the conditions):
equation~\eqref{eq:poly} is independent of $N_{\mathrm{eff}}$ —
fixing $R_\psi = 1$ is a property of S-symmetry alone.

Closes: scale closure conditionally \MC\ (candidate, not proof).\\
Implication for alpha: $\tau = i$ (self-dual point) algebraically motivated
by the S-symmetry structure of $S[\Theta]$ — $B_{\mathrm{phenom}}$ is
motivated but \textbf{NOT DERIVED}.\\
Status: \MC\ — conditional candidate for closing scale closure.
\end{insightbox}

%──────────────────────────────────────────────────────────────────────────────
\section{Implications for C2-iv: Modular Symmetry and the Colour-$\psi$ Coupling}
\label{sec:c2iv}
%──────────────────────────────────────────────────────────────────────────────

\begin{remark}[C2-iv and the $\psi$-circle at $R_\psi = 1$]
\label{rem:c2iv}
If $R_\psi = 1$ is fixed by modular symmetry, then at the self-dual point
$\tau = i$:
\begin{itemize}
  \item $S(i) = i$ — $\tau = i$ is a fixed point.
  \item $\eta(i),\, \vartheta_3(0|i)$ take special values
    (CM points, exactly computable).
  \item The KK spectrum on $S^1_\psi$ is symmetric under $n \to -n$.
\end{itemize}
Question (C2-iv): at the self-dual point $R_\psi = 1$, is colour $\ZZ_3$
naturally coupled to $S^1_\psi$?

Specifically: the rotation $\psi \to \psi + 2\pi/3$ (one third of the
period $2\pi$) is a $\ZZ_3$ symmetry of $S^1_\psi$ with period $2\pi$ at
$R_\psi = 1$. This action commutes with colour $\ZZ_3$ if there exists an
identification
\[
  \psi\text{-rotation by }\tfrac{2\pi}{3}
  \;\longleftrightarrow\;
  \text{colour }\ZZ_3\text{-rotation in SU(3)}.
\]
This is a \emph{candidate closure of C2-iv} \MC\ conditional on modular
fixing of $R_\psi$ and on derivation of colour structure from $\psi$-winding
(C2-i: \OPEN).
\end{remark}

%──────────────────────────────────────────────────────────────────────────────
\section{Numerical Verification}
\label{sec:numerical}
%──────────────────────────────────────────────────────────────────────────────

The script \texttt{tools/modular\_rpsi\_check.py} performs a complete numerical
verification. Results (rounded):
\[
  N_{\mathrm{eff}} = 12,\quad
  C_{\mathrm{unit}} \approx 2.855 \times 10^{-3},\quad
  C_{\mathrm{total}} \approx 3.425 \times 10^{-2}.
\]

\begin{center}
\begin{tabular}{cccc}
  \toprule
  $R_\psi$ & $T_{\mathrm{kin}}$ (S-inv) & $dS_{\mathrm{eff}}/dR$ & $S_{\mathrm{eff}}(R)$ \\
  \midrule
  $0.50$ & 0.02862182 & $-1.4644$ & 0.36395421 \\
  $0.80$ & 0.01745931 & $-0.1412$ & 0.15466341 \\
  $0.90$ & 0.01659830 & $-0.0523$ & 0.14084953 \\
  $\mathbf{1.00}$ & $\mathbf{0.01635533}$ & $\approx 0$ & $\mathbf{0.13701806}$ \\
  $1.10$ & 0.01655402 & $+0.0338$ & 0.14014911 \\
  $1.20$ & 0.01708829 & $+0.0578$ & 0.14866586 \\
  $2.00$ & 0.02862182 & $+0.1734$ & 0.36395421 \\
  \bottomrule
\end{tabular}
\end{center}

\medskip\noindent
S-invariance check: $S_{\mathrm{eff}}(R) = S_{\mathrm{eff}}(1/R)$ numerically
to precision $\sim 10^{-17}$ for all tested values of $R$. \NUM

\medskip\noindent
$N_{\mathrm{eff}}$-independence: $dS/dR|_{R=1} \approx 0$ for $N_{\mathrm{eff}}
\in \{1, 3, 6, 12, 24, 48\}$ (residuals $\lesssim 10^{-16}$, numerical noise). \NUM

\medskip\noindent
Second derivative: $d^2S/dR^2|_{R=1} = 20\,C_{\mathrm{total}} \approx 0.685 > 0$
— confirmed minimum. \NUM

%──────────────────────────────────────────────────────────────────────────────
\section{Summary and Implications}
\label{sec:summary}
%──────────────────────────────────────────────────────────────────────────────

\begin{insightbox}
\textbf{Summary of results (2026-06-11, updated with $Z[\tau]$ chain)}

\begin{enumerate}
  \item \textbf{Polynomial~\eqref{eq:poly}}: S-invariance $+$ stationarity
    gives $R_\psi^6 + R_\psi^4 + R_\psi^2 = 3$, unique positive solution
    $R_\psi = 1$. $N_{\mathrm{eff}}$ cancels. \Lzero\ (algebraic identity).
  \item \textbf{Minimum}: $R_\psi = 1$ is the minimum of $S_{\mathrm{eff}}$,
    $d^2S/dR^2 = 20C > 0$. \Lzero.
  \item \textbf{$Z[\tau_E] = \eta(\tau_E)^{-12}$}: From $\det P = \eta^2$
    \STD\ and $N_{\mathrm{eff}}=12$ \Lone: partition function derived.
    \textbf{\Lone\ conditional} (one-loop approximation).
  \item \textbf{$\psi$ compact}: from SS condition \Lone\
    (Lem.~\ref{lem:SS_BC}) and from periodicity of $Z[\tau]$ \Lzero\
    (Prop.~\ref{prop:psi_compact_from_Z}). Two independent mechanisms.
  \item \textbf{$R_\psi=1$ fixed point of $Z[\tau]$}: $\tau_E=i$ is the
    fixed point of the S-transformation on $Z[\tau_E]$ \Lzero\
    (Thm.~\ref{thm:Rpsi_from_Z}). Agreement with polynomial: robust result
    from two independent paths. \textbf{Overall status: \Lone\ conditional
    (Path 2)}.
  \item \textbf{S-invariance of $S[\Theta]$ in general}: \OPEN\ (condition 1
    for the Casimir path). \emph{Not needed} for Path 2 via $Z[\tau]$.
  \item \textbf{Scale closure EW-1b}: \MC\ — strong candidate.
  \item \textbf{Alpha}: \textbf{NOT DERIVED}. $\tau=i$ algebraically
    preferred via two paths; $B_{\mathrm{phenom}}$ and $\alpha$ remain open.
\end{enumerate}

\begin{center}
\begin{tabular}{lll}
  \toprule
  Path & Result & Status \\
  \midrule
  Casimir polynomial & $R_\psi=1$ minimum of $S_{\mathrm{eff}}$ & \Lzero\ $+$ \MC \\
  $Z[\tau]=\eta^{-12}$ & $R_\psi=1$ fixed point of $Z[\tau]$ & \Lone\ cond. \\
  \bottomrule
\end{tabular}
\end{center}

\textbf{Main result}: $R_\psi = 1$ is derived from the symmetry of the
partition function as \Lone\ conditional (Path 2), with an independent
\Lzero\ confirmation (Path 1). This constitutes a first-principles derivation
of $R_\psi$ in UBT.
\end{insightbox}

\begin{thebibliography}{9}
\bibitem{polchinski_v1}
  J.~Polchinski, \textit{String Theory, Vol.~1}.
  Cambridge University Press, 1998.

\bibitem{ambjorn_wolfram_1983}
  J.~Ambj\o{}rn and S.~Wolfram,
  ``Properties of the vacuum. 1. Mechanical and thermodynamic,''
  \textit{Ann.\ Phys.}\ \textbf{147}, 1 (1983).

\bibitem{apostol_modular}
  T.~M.~Apostol, \textit{Modular Functions and Dirichlet Series in
  Number Theory}, 2nd~ed.\ Springer, 1990.

\bibitem{su2_twist}
  I.~D.~Jaroš, ``SU(2) Scherk--Schwarz Twist and $N_{\mathrm{eff}}=12$
  in UBT,'' \texttt{research\_tracks/quantum\_ubt/su2\_twist\_neff12.tex},
  2026.
\end{thebibliography}

\end{document}
\fi
